\documentclass[11pt]{article}
\usepackage{amsmath,amssymb}
\usepackage[margin=1in]{geometry}

\title{Orbital Rotation Gradient and Hessian for a General JK-only Functional}
\author{}
\date{}

\begin{document}
\maketitle

\section*{Setup}

Fixed natural occupations $n_p$, $D_{pq}=n_p\delta_{pq}$. A JK-only two-electron
energy functional with a general symmetric occupation-coupling function
$f(n_a,n_b)=f(n_b,n_a)$, $f(n,n)=n$, replacing the two-particle density matrix
elements as
\begin{equation}
\Gamma_{pqpq} = \tfrac12\,n_p n_q \quad (\text{Hartree/direct}), \qquad
\Gamma_{pqqp} = -\tfrac12\,f(n_p,n_q) \quad (\text{exchange}),
\end{equation}
so that
\begin{equation}
E_2 = \tfrac12\sum_{pq} n_p n_q\,\langle pq|pq\rangle
\;-\; \tfrac12\sum_{pq} f(n_p,n_q)\,\langle pq|qp\rangle .
\end{equation}

The M\"uller functional is the special case $f(n_a,n_b)=\sqrt{n_a n_b}$.

\bigskip
\noindent\textbf{Important:} this $\Gamma$ does \emph{not} satisfy the
fermionic antisymmetry $\Gamma_{pqrs}=-\Gamma_{qprs}$ of a genuine two-particle
density matrix (it is an $N$-representability-violating reconstruction, valid
only in the idempotent/integer-occupation limit). The gradient and Hessian
below are therefore derived directly from the explicit energy expression
above (differentiating each integral leg under orbital rotation), \emph{not}
from the operator-commutator machinery used for the general $\Gamma$ case,
since that machinery relied on antisymmetry at two separate steps.

\section*{Gradient}

Define
\begin{equation}
V_{qp} \;\equiv\; h_{qp} + \sum_t n_t\,\langle tq|tp\rangle .
\end{equation}

\begin{equation}
\boxed{\,
g_{pq} \;=\; (n_q-n_p)\,V_{qp}
\;+\; \sum_t \big[f(n_p,n_t)-f(n_q,n_t)\big]\,\langle qt|tp\rangle
\,}
\end{equation}

\section*{Hessian}

\begin{align}
G_{pq,rs} \;=\;\;
& (n_q-n_p)\Big[-\delta_{qr}\,V_{sp} + \delta_{ps}\,V_{qr}
+ (n_s-n_r)\,\langle sq|rp\rangle\Big] \notag\\[4pt]
& -\,\delta_{qr}\sum_t \big[f(n_p,n_t)-f(n_r,n_t)\big]\,\langle st|tp\rangle
\;+\; \delta_{ps}\sum_t \big[f(n_s,n_t)-f(n_q,n_t)\big]\,\langle qt|tr\rangle \notag\\[4pt]
& +\,\big[f(n_p,n_s)-f(n_p,n_r)-f(n_q,n_s)+f(n_q,n_r)\big]\,\langle qs|rp\rangle
\end{align}

\bigskip
\noindent\textbf{Special case (M\"uller):} set $f(n_a,n_b)=\sqrt{n_an_b}$
everywhere above.

\bigskip
\noindent\textbf{Status note.} This derivation does not rely on 2-RDM
antisymmetry (appropriately, since the JK-only $\Gamma$ lacks it), but it has
\emph{not} been independently cross-checked by a second derivation route or a
symbolic tool. Given the nonlinear occupation-coupling $f$ and the amount of
index bookkeeping, a numerical finite-difference validation of
$E(\boldsymbol\kappa)$ against this gradient and Hessian (small system,
generic fractional occupations, both for $f=\sqrt{\phantom{x}}$ and at least
one other $f$) is strongly recommended before use in production code.

\end{document}
